Un tore à section rectangulaire porte $N = \v{N}$ spires parcourues par un courant de \qty{\v{I}}{A}. Son rayon intérieur est $a = \qty{\v{a}}{cm}$, son rayon extérieur $b = \qty{\v{b}}{cm}$ et sa hauteur $h = \qty{\v{h}}{cm}$. ($\mu_0 = 4\pi \times \qty{e-7}{T.m/A}$)

\begin{enumerate}[label=\alph*)]
    \item À l'aide du théorème d'Ampère, déterminez le champ magnétique à l'intérieur du tore, au rayon intérieur $a$.
    \item Déterminez le champ magnétique au rayon extérieur $b$.
    \item Le champ n'est pas uniforme dans la section. Déterminez le flux magnétique à travers une seule spire.
\end{enumerate}
